% ECONOMICS_PROBLEM: {"id": "E32-428", "title": "Structural change in cyclical models", "jel_code": "E32", "source_id": "OP-0240"}
% !TeX program = xelatex
\documentclass[11pt,a4paper]{article}
\usepackage[margin=23mm]{geometry}
\usepackage{fontspec}
\setmainfont{Latin Modern Roman}
\usepackage{amsmath,amssymb,mathtools}
\usepackage{xcolor,enumitem,booktabs,longtable,array,xurl,hyperref}
\definecolor{muted}{HTML}{53636C}
\hypersetup{colorlinks=true,urlcolor=blue,linkcolor=blue}
\setlength{\parindent}{0pt}
\setlength{\parskip}{5pt}
\newcommand{\R}{\mathbb R}
\newcommand{\E}{\mathbb E}
\newcommand{\N}{\mathbb N}
\newcommand{\ind}{\mathbf 1}
\newcommand{\Var}{\operatorname{Var}}
\newcommand{\Cov}{\operatorname{Cov}}
\newcommand{\supp}{\operatorname{supp}}
\DeclareMathOperator*{\argmin}{arg\,min}
\DeclareMathOperator*{\argmax}{arg\,max}
\newcommand{\entrymeta}[3]{{\sffamily\small\color{muted}\textbf{\texttt{#1}}\quad|\quad #2\par #3\par}}
\newcommand{\Problemref}[1]{\texttt{#1}}
\newcommand{\JELref}[2]{\texttt{#2} (JEL \texttt{#1})}
\begin{document}
\section*{Structural change in cyclical models}
\textbf{Problem ID:} \texttt{E32-428}\quad \textbf{JEL:} \texttt{E32}\par
% BEGIN ECONOMICS STATEMENT
\label{problem:OP-0240}
{\sffamily\small\color{muted}Original subject group: Business cycles and macro dynamics\par}
\label{definitions:problem:30007521}
\textbf{Audit classification:} Published research agenda. Structural model development is listed; the joint slow-state intervention target is editorial. Verification concerns the cited version, not first priority or proof that this exact editorial specification remains unsolved on October 8, 2026.
\subsection*{Definitions and precise research target}
\paragraph{Editorial mathematical specification.} Fix a baseline macroeconomic model with equilibrium state \(s_t\), policy rule \(a_t\), and observable vector \(z_t\). Add measurable slow states \(v_t=(v_t^C,v_t^D,v_t^N)\), denoting climate exposure, digital-capital penetration, and demographic composition. Their definitions, units, data vintages, and transition laws must be specified. Let \(M_0\) hold those states fixed at a benchmark and let \(M_1\) allow them to affect production, demand, labor supply, and financial constraints through an identified finite-dimensional parameter \(\theta\).
For a common identified monetary or sectoral intervention \(a\) and benchmark action \(a^0\), define
\[
\begin{split}
\Delta(h,v)={}&\big(E_{M_1}[z_{t+h}\mid\operatorname{do}(a_t=a),v_t=v]
-E_{M_1}[z_{t+h}\mid\operatorname{do}(a_t=a^0),v_t=v]\big)\\
&-\big(E_{M_0}[z_{t+h}\mid\operatorname{do}(a_t=a)]
-E_{M_0}[z_{t+h}\mid\operatorname{do}(a_t=a^0)]\big),
\end{split}
\]
Determine which channels and interactions among the three slow states are required to predict \(\Delta(h,v)\) over a declared horizon and range of environments. Identify causal variation separately from concurrent trends, and avoid treating every slow movement as an independent shock. Compare nested restrictions using held-out economies, policy episodes, and forecast distributions; estimate the uncertainty attributable to each added channel. A successful result is a tractable, empirically supported extension with improved policy-response predictions under specified structural transitions. It is not a demand for one universal model covering all possible climate, technological, and demographic futures.

Specify \(s_{t+1}=F_\theta(s_t,v_t,a_t,\varepsilon_{t+1})\), \(v_{t+1}=G_\theta(v_t,\zeta_{t+1})\), and \(z_t=H_\theta(s_t,v_t)\), with fixed shock laws and initial state. Define \(M_0\) by clamping \(v\) to a stated value in these same equations, not by separately recalibrating all other parameters. The intervention replaces one decision-rule equation by the assigned action \(a_t=a\); later actions follow the given common rule. Compare response contrasts relative to each model's no-intervention baseline to avoid confusing a level difference with altered transmission. Climate exposure may be a physical state, digitalization a capital stock, and demography a distribution; their units and state transitions cannot be treated as interchangeable.
\subsection*{Variants and assumption changes}
The following are \emph{editorial variants}, not additional published open conjectures.
\begin{enumerate}
\item \emph{One-state restriction.} Vary only measured climate exposure and hold digital capital and demographic composition fixed. Identify the difference in monetary impulse responses between two exposure values using a stated interaction instrument.
\item \emph{Joint nonlinear transition.} Permit endogenous digital investment and migration to respond to climate exposure. Compare the full response with the sum of the three single-state responses, defining the residual as an interaction rather than an unexplained independent shock.
\end{enumerate}
\subsection*{References and source audit}
\begin{itemize}
\item European Central Bank. \emph{Forecasting and business cycle analysis}. Undated. Locator: Model development, selected areas of focus. Primary text checked. \url{https://www.ecb.europa.eu/pub/economic-research/research_agenda/forecasting/html/index.en.html}.
\end{itemize}
Structural model development is listed; the joint slow-state intervention target is editorial.
\begin{enumerate}\raggedright
\item\hypertarget{definitions:source:30007521:1}{} European Central Bank. Undated / date unverified. \emph{Forecasting and business cycle analysis}. Model development, selected areas of focus.
\url{https://www.ecb.europa.eu/pub/economic-research/research_agenda/forecasting/html/index.en.html}.
\end{enumerate}

\subsection*{Common conventions: Source mathematical conventions}

These conventions apply to each entry unless explicitly replaced.

\textbf{Models and identification.} A primitive model \(m\) specifies all probability laws, feasible choices, information, equilibrium and measurement rules used by its target. The admissible class \(\mathcal M\) is fixed independently of outcomes. If \(P_O(m)\) is its observable probability law and \(\tau(m)\) the target, its identified set at observable law \(P\) is
\[
 \mathcal I_\tau(P)=\{\tau(m):m\in\mathcal M,\ P_O(m)=P\}.
\]
Point identification means that this set is a singleton. An empty set signals incompatibility of data and assumptions. Coordinate bounds need not describe the joint set. Bounds are sharp only if every retained target is attainable or arbitrarily approachable by compatible models. Finite bounds require explicit restrictions on unobserved quantities; unrestricted confounding, hidden wealth, extrapolation or arbitrary equilibrium selection can yield uninformative sets.

\textbf{Causal interventions.} A potential outcome \(Y_i(p)\) is the outcome under a complete policy regime \(p\), including timing, financing, information and market spillovers. The target population and units are fixed before treatment. With interference, use the complete assignment vector or a justified exposure mapping; individual no-interference notation alone cannot identify an economy-wide intervention. A difference of mean potential outcomes does not require the cross-world joint distribution. Conditioning on post-treatment migration, survival, enrollment or employment changes the estimand and needs separate justification.

\textbf{Statistical statements.} An estimator, confidence set, forecast or test specifies a sampling law, model class and loss/coverage criterion. Uniform \((1-\alpha)\)-coverage means \(\inf_{m\in\mathcal M}\Pr_m\{\tau(m)\in C_n\}\geq1-\alpha\), or an explicitly asymptotic version. The whole parameter space trivially has coverage, so informativeness requires a stated width, risk or power objective. A forecasting improvement is relative to a named benchmark, information set and evaluation population.

\textbf{Equilibrium and optimization.} An equilibrium comprises feasible plans, optimality, beliefs where needed, budget constraints and all market/resource clearing. The primitives or a stated benchmark define each correspondence \(\mathcal E(p;m)\). Nonemptiness is assumed or proved, never inferred just from compact policy sets. A selected-equilibrium target uses a declared selection rule; a robust target quantifies over all permitted equilibria. Use suprema or infima unless attainment is established. An \(\varepsilon\)-optimal policy is feasible and within \(\varepsilon\) of the finite optimum value. Report an empty feasible set separately.

\textbf{Welfare and accounting.} Normative utility, interpersonal normalization, social weights, numeraire, discounting and the use of public receipts are primitives. Behavioral utility need not coincide with the welfare criterion. Household welfare includes ownership income and rents once; adding profits again to compensating variation generally double-counts them. A monetary equivalent is defined by an explicit transfer and price convention, with monotonicity and range conditions. Average effects, mechanism contributions, transfers and welfare losses are different targets. Infinite sums require convergence and dynamic feasible plans require terminal or no-Ponzi conditions.

\textbf{Source versions.} A working-paper date, revision date and journal publication year are distinguished. A DOI for a journal version is not automatically the DOI of an earlier manuscript. Locators refer to the stated version and printed pages where specified. A metadata-only check does not verify a claim about a full-text conclusion. Every entry's source subsection narrows the provenance claim accordingly.
% END ECONOMICS STATEMENT
\end{document}
